% Standalone resolution manuscript generated from solution.md.
% Compile with XeLaTeX; author and review metadata are maintained in Markdown.
\documentclass[10pt,a4paper]{article}
\usepackage[margin=24mm,headheight=15pt]{geometry}
\usepackage{fontspec}
\setmainfont{DejaVuSerif.ttf}[BoldFont=DejaVuSerif-Bold.ttf,ItalicFont=DejaVuSerif-Italic.ttf,BoldItalicFont=DejaVuSerif-BoldItalic.ttf]
\setsansfont{DejaVuSans.ttf}[BoldFont=DejaVuSans-Bold.ttf,ItalicFont=DejaVuSans-Oblique.ttf]
\setmonofont{DejaVuSansMono.ttf}
\usepackage{amsmath,amssymb,unicode-math}
\setmathfont{latinmodern-math.otf}
\usepackage{xcolor,microtype,parskip,needspace,fancyhdr}
\usepackage{bookmark}
\usepackage{xurl}
\hypersetup{colorlinks=true,urlcolor=blue,linkcolor=blue,pdftitle={MI-28:
a determinant inequality for all nonnegative base
powers},pdfauthor={George Stepaniants},pdflang={en-GB}}
\urlstyle{same}
\setlength{\emergencystretch}{3em}
\providecommand{\tightlist}{\setlength{\itemsep}{0pt}\setlength{\parskip}{0pt}}
\providecommand{\pandocbounded}[1]{#1}
\setcounter{secnumdepth}{0}
\allowdisplaybreaks[1]
\widowpenalty=10000
\clubpenalty=10000
\pagestyle{fancy}
\fancyhf{}
\fancyhead[L]{\small\sffamily Verified resolution / MI-28}
\fancyhead[R]{\small\sffamily 11 September 2026}
\fancyfoot[L]{\parbox[b]{0.9\textwidth}{\fontsize{8}{10}\selectfont Independent
Codex-agent review; not external human peer review or formal
certification.}}
\fancyfoot[R]{\footnotesize\thepage}
\begin{document}
{\raggedright\Large\bfseries MI-28: a determinant inequality for all
nonnegative base powers\par}
\vspace{0.4em}
{\normalsize\bfseries George Stepaniants\par}
{\small Department of Computing and Mathematical Sciences, California
Institute of Technology, Pasadena, California, USA\par}
{\small\href{mailto:gstepan@caltech.edu}{gstepan@caltech.edu}\par}
\vspace{0.6em}
\pagestyle{plain}

For positive definite matrices \(A,B\), every \(k\ge0\), and
\(0\le p\le2\), we prove

\[
\det(A^k+|AB|^p)\ge\det(A^k+A^pB^p).
\]

In fact, a log-majorization of the corresponding normalized matrices
holds. The previously proved range \(k\ge2\) is due to Ghabries, Abbas,
Mourad, and Assi. We treat the remaining base powers using a normalized
order implication. Two applications of the Furuta inequality establish
complementary parameter regions. The identity
\(\lvert |AB|^{-1}B\rvert=A^{-1}\) interchanges the two parameters and
closes the remaining range.

\textbf{Verification status: independent-agent PASS.} Developed with
substantial ChatGPT/Codex assistance. A separate agent audited the
complete analytic argument, its parameter ranges and its primary-source
applications.
\href{https://github.com/ajt60gaibb/OpenProblemsInNLA/blob/main/references/stepaniants-mi28-2026-09-11/verification/reviews/MI-28-review.md}{Detailed
independent review}. This is independent agent review, not external
human peer review or formal verification.

\subsection{Statement and notation}\label{statement-and-notation}

Throughout, \(A,B\in\mathbb C^{n\times n}\) are positive definite,
matrix powers are defined by spectral functional calculus, and
\(\|\cdot\|\) denotes the operator norm. For a matrix \(M\), write
\(|M|=(M^*M)^{1/2}\). Eigenvalues of a positive definite matrix are
listed in decreasing order. For positive vectors \(u,v\in(0,\infty)^n\),
the notation \(u\prec_{\log}v\) means

\[
\prod_{i=1}^j u_i\le\prod_{i=1}^j v_i\quad(1\le j<n),
\qquad
\prod_{i=1}^n u_i=\prod_{i=1}^n v_i.
\]

\subsubsection{Theorem 1 --- The full determinant and log-majorization
inequalities}\label{theorem-1-the-full-determinant-and-log-majorization-inequalities}

For every \(A,B>0\), \(k\ge0\), and \(0\le p\le2\),

\[
\lambda\!\left(A^{(p-k)/2}B^pA^{(p-k)/2}\right)
\prec_{\log}
\lambda\!\left(A^{-k/2}|AB|^pA^{-k/2}\right).
\]

Consequently,

\[
\det(A^k+|AB|^p)\ge\det(A^k+A^pB^p).
\]

The determinant inequality is the positive definite formulation of
\href{https://github.com/ajt60gaibb/OpenProblemsInNLA/blob/main/matrix-inequalities-and-norms/MI-28/README.md}{MI-28},
originating in
\href{https://www.researchgate.net/publication/361793582_Contributions_to_Matrix_Inequalities_and_Some_Applications}{Ghabries's
thesis}, §2.5 and final Open Problems, Problem 1, pp.~109--110. Although
\(A^k+A^pB^p\) need not be Hermitian, its determinant is positive: after
factoring out \(A^k\), the remaining product \(A^{p-k}B^p\) is similar
to the positive definite matrix on the left of the log-majorization. The
range \(k\ge2\), \(0\le p\le2\) follows already from
\href{https://www.researchgate.net/publication/342908148_A_proof_of_a_conjectured_determinantal_inequality}{Ghabries--Abbas--Mourad--Assi,
Lemma 2.5}; the precise substitution appears in the proof of Theorem 1
below.

\subsection{Two order implications}\label{two-order-implications}

We use the Löwner--Heinz inequality

\[
0<X\le Y,\quad 0\le t\le1
\quad\Longrightarrow\quad X^t\le Y^t,
\]

and the following convention for the Furuta inequality: if \(X\ge Y>0\),
\(r\ge0\), \(a\ge0\), \(q\ge1\), and \((1+r)q\ge a+r\), then

\[
(X^{r/2}Y^aX^{r/2})^{1/q}\le X^{(a+r)/q}.
\]

Furuta's original paper is listed in the references. Both facts are also
restated in
\href{https://www.researchgate.net/publication/255605812_The_Furuta_inequality_with_negative_powers}{Tanahashi,
Propositions 1--2}. The \(r\) used here is twice the sandwich exponent
in that reference.

For \(k,p>0\), let \(\mathcal I(k,p)\) denote the assertion that, for
every \(A,B>0\), with \(D=|AB|\),

\[
D^p\le A^k\quad\Longrightarrow\quad B^p\le A^{k-p}.
\]

This notation describes a universal order implication. It makes no
assumption that \(A\) and \(B\) commute.

\subsubsection{Lemma 2 --- The lower-power Furuta
implication}\label{lemma-2-the-lower-power-furuta-implication}

If \(k>0\) and \(k/(k+1)\le p\le1\), then \(\mathcal I(k,p)\) holds.

\textbf{Proof.} Assume \(D^p\le A^k\). In the Furuta inequality, choose

\[
X=A^k,\qquad Y=D^p,\qquad r=\frac2k,\qquad
a=\frac2p,\qquad q=2.
\]

The condition \((1+r)q\ge a+r\) is exactly \(p\ge k/(k+1)\). Thus

\[
(AD^2A)^{1/2}\le A^{1+k/p}.
\]

Since \(D^2=BA^2B\), we have \(AD^2A=(ABA)^2\), and its positive square
root is \(ABA\). Congruence by \(A^{-1}\) gives \(B\le A^{k/p-1}\).
Applying Löwner--Heinz with exponent \(p\in(0,1]\) proves
\(B^p\le A^{k-p}\). \(\square\)

\subsubsection{Lemma 3 --- The higher-power Furuta
implication}\label{lemma-3-the-higher-power-furuta-implication}

If \(k>0\) and \(1\le p\le2\), then \(\mathcal I(k,p)\) holds.

\textbf{Proof.} Assume \(D^p\le A^k\), and set

\[
C=ABA,\qquad s=\frac{p(k+2)}{k+p}.
\]

The hypotheses give \(1\le p\le s\le2\). Apply the Furuta inequality
with

\[
X=A^k,\qquad Y=D^p,\qquad r=\frac2k,\qquad
a=\frac2p,\qquad q=\frac2s.
\]

Here \(q\ge1\) and \((1+r)q=a+r\). Using \(AD^2A=C^2\), we obtain

\[
C^s=(AD^2A)^{s/2}\le A^{k+2}.
\]

Because \(0<2/(k+2)\le1\), Löwner--Heinz followed by inversion gives

\[
A^{-2}\le C^{-2s/(k+2)}=C^{-2p/(k+p)}.
\]

For every invertible \(S\) and real \(p\), a singular-value
decomposition gives

\[
(SS^*)^p=S(S^*S)^{p-1}S^*.
\]

Apply this identity to \(S=A^{-1}C^{1/2}\), for which \(SS^*=B\):

\[
B^p=A^{-1}C^{1/2}
(C^{1/2}A^{-2}C^{1/2})^{p-1}
C^{1/2}A^{-1}.
\]

Congruence in the inequality for \(A^{-2}\), followed by Löwner--Heinz
with \(0\le p-1\le1\), therefore yields

\[
\begin{aligned}
B^p
&\le A^{-1}C^{\,1+(p-1)(k-p)/(k+p)}A^{-1}\\
&=A^{-1}C^{\,p(k-p+2)/(k+p)}A^{-1}.
\end{aligned}
\]

The exponent of \(C\) in the last line equals \(s\theta\), where

\[
\theta=\frac{k-p+2}{k+2}\in[0,1].
\]

Applying Löwner--Heinz to \(C^s\le A^{k+2}\) gives
\(C^{s\theta}\le A^{k-p+2}\). Substitute this in the bound for \(B^p\)
to conclude \(B^p\le A^{k-p}\). \(\square\)

\textbf{Remark 4 --- Relation to the negative-power Furuta inequality.}
The argument after \(C^s\le A^{k+2}\) is a direct proof of the needed
specialization of the negative-power Furuta inequality. In the
convention used here, the relevant parameters in
\href{https://www.researchgate.net/publication/255605812_The_Furuta_inequality_with_negative_powers}{Tanahashi,
Theorem 3} are

\[
X=A^{k+2},\quad Y=C^s,\quad
a=\frac{k+p}{p(k+2)},\quad q=\frac1p,\quad r=-\frac2{k+2}.
\]

The displayed proof verifies that specialization using only
Löwner--Heinz and the singular-value identity.

\subsection{Interchanging the
parameters}\label{interchanging-the-parameters}

\subsubsection{Lemma 5 --- The parameter
swap}\label{lemma-5-the-parameter-swap}

Let \(0<p\le k\). If \(\mathcal I(p,k)\) holds, then \(\mathcal I(k,p)\)
holds.

\textbf{Proof.} Assume \(D^p\le A^k\), where \(D=|AB|\), and define

\[
\widetilde A=D^{-1},\qquad \widetilde B=B.
\]

The order of the factors is essential. Directly,

\[
|\widetilde A\widetilde B|^2
=BD^{-2}B
=B(B^{-1}A^{-2}B^{-1})B
=A^{-2},
\]

so \(|\widetilde A\widetilde B|=A^{-1}\). Inverting the assumed order
gives

\[
|\widetilde A\widetilde B|^k=A^{-k}
\le D^{-p}=\widetilde A^p.
\]

Apply \(\mathcal I(p,k)\) to this pair:

\[
B^k\le\widetilde A^{p-k}=D^{k-p}.
\]

Since \(0<p/k\le1\), Löwner--Heinz applied to this inequality gives

\[
B^p\le D^{p(k-p)/k}.
\]

Since \(0\le(k-p)/k\le1\), applying it also to \(D^p\le A^k\) gives

\[
D^{p(k-p)/k}\le A^{k-p}.
\]

Together these inequalities prove \(\mathcal I(k,p)\). If \(p=k\), the
second exponent is zero and its inequality is simply \(I\le I\); the
argument still applies. \(\square\)

\subsubsection{Corollary 6 --- All base powers up to
two}\label{corollary-6-all-base-powers-up-to-two}

For \(0<k\le2\) and \(0<p\le2\), the implication \(\mathcal I(k,p)\)
holds.

\textbf{Proof.} If \(p\ge1\), use Lemma 3. Suppose \(0<p<1\). If
\(p\ge k/(k+1)\), use Lemma 2. In the remaining case,

\[
p<\frac{k}{k+1}<k\le2.
\]

We establish the swapped implication \(\mathcal I(p,k)\). When
\(k\le1\), its exponent \(k\) satisfies

\[
\frac{p}{p+1}<p<k\le1,
\]

so Lemma 2, with the parameters interchanged, applies. When
\(1\le k\le2\), Lemma 3, with the parameters interchanged, applies.
Lemma 5 now proves \(\mathcal I(k,p)\). \(\square\)

\subsection{Log-majorization and the
determinant}\label{log-majorization-and-the-determinant}

\textbf{Proof of Theorem 1.} First take \(0<k\le2\), \(0<p\le2\), and
write

\[
H=A^{(p-k)/2}B^pA^{(p-k)/2},\qquad
Z=A^{-k/2}D^pA^{-k/2}.
\]

Both matrices are homogeneous of degree \(p\) in \(B\). Let
\(c=\|Z\|>0\), and replace \(B\) by \(c^{-1/p}B\). The resulting \(Z\)
is at most \(I\), equivalently \(D^p\le A^k\) for the scaled pair.
Corollary 6 gives \(B^p\le A^{k-p}\). Congruence by \(A^{(p-k)/2}\)
shows that the scaled \(H\) is at most \(I\). Scaling back proves

\[
\|H\|\le\|Z\|.
\]

For every \(1\le j\le n\), apply this norm inequality to
\(\bigwedge^j A\) and \(\bigwedge^j B\). These exterior powers are
positive definite and preserve products, inverses, and spectral powers.
In particular,

\[
\left|(\bigwedge\nolimits^j A)(\bigwedge\nolimits^j B)\right|
=\bigwedge\nolimits^j |AB|.
\]

The operator norm of an exterior power of a positive definite matrix is
the product of its \(j\) largest eigenvalues. We thus obtain all
partial-product inequalities in the claimed log-majorization. The full
determinants agree:

\[
\det H=\det Z=\det(A)^{p-k}\det(B)^p.
\]

This proves the log-majorization for \(0<k\le2\), \(0<p\le2\).

For \(k\ge2\), use the previously established result
\href{https://www.researchgate.net/publication/342908148_A_proof_of_a_conjectured_determinantal_inequality}{Ghabries--Abbas--Mourad--Assi,
Lemma 2.5}. In its notation, for \(0\le t\le s\le K\),

\[
\lambda(X^{K-t}Y^t)
\prec_{\mathrm{wlog}}
\lambda\!\left(
X^{K/2}(Y^{s/2}X^{-s}Y^{s/2})^{t/s}X^{K/2}\right).
\]

Take \(X=A^{-1}\), \(Y=B\), \(K=k\), \(s=2\), and \(t=p\). The two
eigenvalue lists are exactly those of \(H\) and \(Z\); the product
\(A^{p-k}B^p\) is similar to \(H\). Equality of the full determinants
changes this weak log-majorization into the claimed log-majorization.
This invokes the original range \(0\le t\le s\le K\) without altering
any of its parameter restrictions.

When \(p=0\), the two matrices in the claimed log-majorization are both
\(A^{-k}\), so equality holds. For \(k=0\), \(p>0\), let
\(k\downarrow0\) in the established range \(0<k\le2\). Functional
calculus and the ordered eigenvalues are continuous on the positive
definite cone, and each partial-product inequality and the equality of
the full determinants persists. Thus the log-majorization holds at every
stated endpoint.

Finally, log-majorization is ordinary majorization of the logarithms of
the eigenvalues. Applying the convex function \(t\mapsto\log(1+e^t)\)
gives

\[
\det(I+H)\le\det(I+Z).
\]

Indeed, the corresponding inequality for the sums of this convex
function is the standard convex-function characterization of
majorization. Since \(A^{p-k}B^p\) is similar to \(H\),

\[
\begin{aligned}
\det(A^k+A^pB^p)&=\det(A^k)\det(I+H),\\
\det(A^k+D^p)&=\det(A^k)\det(I+Z).
\end{aligned}
\]

Multiplying by the positive number \(\det(A^k)\) proves the determinant
inequality. \(\square\)

\textbf{Remark 7 --- Scope.} The theorem treats every dimension and
every parameter pair in the canonical positive definite statement of
MI-28. For \(k>0\) and \(p>0\), the positive semidefinite extension
follows by applying the determinant inequality to \(A+\varepsilon I\)
and \(B+\varepsilon I\), then letting \(\varepsilon\downarrow0\). The
positive definite formulation avoids conventions for zeroth powers of
singular matrices.

\subsection{References}\label{references}

\begin{enumerate}
\def\labelenumi{\arabic{enumi}.}
\item
  \textbf{MI-28.} ajt60gaibb/OpenProblemsInNLA. \emph{Ghabries's
  determinant comparison for arbitrary nonnegative base powers}.
  \href{https://github.com/ajt60gaibb/OpenProblemsInNLA/blob/main/matrix-inequalities-and-norms/MI-28/README.md}{Canonical
  problem statement}.
\item
  \textbf{Ghabries's thesis.} M. M. Ghabries. \emph{Contributions to
  Matrix Inequalities and Some Applications}. PhD thesis, University of
  Angers and Lebanese University, 2022. Section 2.5 and final Open
  Problems, Problem 1, pp.~109--110.
  \href{https://www.researchgate.net/publication/361793582_Contributions_to_Matrix_Inequalities_and_Some_Applications}{Author-uploaded
  thesis}.
\item
  \textbf{Ghabries--Abbas--Mourad--Assi (2020).} M. M. Ghabries, H.
  Abbas, B. Mourad, and A. Assi. \emph{A proof of a conjectured
  determinantal inequality}. Linear Algebra and its Applications
  \textbf{605} (2020), 21--28. Lemma 2.5 and Theorem 1.1.
  \href{https://doi.org/10.1016/j.laa.2020.07.013}{DOI:
  10.1016/j.laa.2020.07.013};
  \href{https://www.researchgate.net/publication/342908148_A_proof_of_a_conjectured_determinantal_inequality}{author-uploaded
  full text}.
\item
  \textbf{Furuta (1987).} T. Furuta. \emph{\(A\ge B\ge O\) assures
  \((B^rA^pB^r)^{1/q}\ge B^{(p+2r)/q}\) for \(r\ge0\), \(p\ge0\),
  \(q\ge1\) with \((1+2r)q\ge p+2r\)}. Proceedings of the American
  Mathematical Society \textbf{101} (1987), 85--88.
\item
  \textbf{Tanahashi (1999).} K. Tanahashi. \emph{The Furuta inequality
  with negative powers}. Proceedings of the American Mathematical
  Society \textbf{127} (1999), 1683--1692. Propositions 1--2 and Theorem
  3. \href{https://doi.org/10.1090/S0002-9939-99-04705-X}{DOI:
  10.1090/S0002-9939-99-04705-X};
  \href{https://www.researchgate.net/publication/255605812_The_Furuta_inequality_with_negative_powers}{author-uploaded
  full text}.
\end{enumerate}
\end{document}
